\section{新推导：保结构海水涡流 Galerkin 与同源端口阻抗}\label{sec:galerkin}
\subsection{问题范围及导电电流函数空间}
本节\textbf{首次将已有 BFZI 解析场源与新的保结构海水电流低阶模型相结合}。
限定如下：谐波峰值相量 \(e^{\jj\omega t}\)、磁准静态、线性互易介质、
均匀 \(\mu=\mu_0\)、海水电导率 \(0<\sigma_{\min}\leq\sigma_s(\bm x)\leq\sigma_{\max}<\infty\)、
绝缘封装及海水外部的正确衰减条件，端口回路与海水界面不相交。
所有基函数和积分取实值几何/材料，复数仅由谐波场系数引入。
静态位姿可任意合法改变，但暂不纳入运动感应电动势。
对可计算有限域 \(\Omega_s\)（或合适的有限能量无限域空间），定义
\begin{equation}
\mathcal V(\Omega_s)
=\{\bm v\in H(\mathrm{div};\Omega_s):
\nabla\cdot\bm v=0,\ 
\bm n\cdot\bm v|_{\partial\Omega_s}=0\quad
\text{（有限域中含人工绝缘外边界）}\}.
\label{eq:sea-solenoidal-space}
\end{equation}
对无限海水，还需要无穷远处衰减、能量有限及适当的函数空间闭合；
后文所有有限维证明先在边界良定、积分收敛的条件下成立。
\textbf{多连通域中仅使用局部 curl 基可能遗漏拓扑环流自由度}；
基空间必须含必要的全局回路模态或采用等价的连续性投影。
这是一项待落实的具体离散构造，不可把其存在性等同于已完成快速实现。

令 \(\bm W_m(\bm x;\mathcal G)\in\mathcal V\) 为 \(r\) 个线性独立的
实向量基，封装与位姿以 \(\mathcal G\) 表示；对任意端口电流列
\(\bm I\in\mathbb C^P\)，采用
\begin{equation}
\bm J_s(\bm x)=\sum_{m=1}^r z_m\bm W_m(\bm x)
=:\bm W(\bm x)\bm z .
\label{eq:sea-reduced-current}
\end{equation}
\(\bm W\) 的列代表电流密度模式，\(\bm z\) 为复数电流模式系数。
允许 \(\bm W\) 随几何变化，但在每一次线性频域求解中应作为固定空间；
随温度改变的 \(\sigma_s\) 进入下面的 \(\bm R_s\)。

\subsection{从 \(A\)--\(\phi\) 方程得到约束弱式}
定义自由空间积分核
\begin{equation}
(\mathcal G_0\bm f)(\bm x)
=\int_{\Omega_s}\frac{\bm f(\bm y)}
{4\pi\|\bm x-\bm y\|}\,\mathrm d\bm y,
\label{eq:free-g}
\end{equation}
以及由 BFZI 完整端口路径（含明确馈/回流）形成的单位端口线源
\(\bm C_p(\bm x)=\int_{\Gamma_p}(4\pi\|\bm x-\bm r'\|)^{-1}
\,\mathrm d\bm\ell'\)。
在 \(\mu=\mu_0\) 条件下，
\[
\bm A(\bm x)
=\mu_0\sum_{p=1}^P \bm C_p(\bm x)I_p
+\mu_0\mathcal G_0\bm J_s(\bm x),\quad
\sigma_s^{-1}\bm J_s=-\jj\omega\bm A-\nabla\phi .
\]
取任意 \(\bm v\in\mathcal V\) 测试，上式与 \(\bm v\) 做
\(L^2\) 内积并对 \(\Omega_s\) 积分。
对采用人工绝缘外边界的\emph{有限域}，
由分部积分、\(\nabla\cdot\bm v=0\) 和\emph{整个边界}
的测试场零法向迹，可得
\(\int_{\Omega_s}\bm v\cdot\nabla\phi\,\mathrm dV=0\)。
若人工外边界允许非零法向通量，则该边界项
不可直接消去；详见第 \ref{sec:g1-foundation-recheck} 节。
故有
\begin{equation}
\int_{\Omega_s}\sigma_s^{-1}\bm v\cdot\bm J_s\,\mathrm dV
+\jj\omega\mu_0\int_{\Omega_s}\bm v\cdot\mathcal G_0\bm J_s\,\mathrm dV
=-\jj\omega\mu_0\sum_p I_p\int_{\Omega_s}\bm v\cdot\bm C_p\,\mathrm dV .
\label{eq:sea-weak}
\end{equation}
采用 \(\bm v=\bm W_m\) 形成有限维矩阵
\begin{align}
[\bm R_s]_{mn}
&=\int_{\Omega_s}\sigma_s^{-1}\bm W_m\cdot\bm W_n\,\mathrm dV,
\label{eq:Rs}\\
[\bm L_s]_{mn}
&=\mu_0\int_{\Omega_s}\bm W_m\cdot\mathcal G_0\bm W_n\,\mathrm dV,
\label{eq:Ls}\\
[\bm K]_{mp}
&=\mu_0\int_{\Omega_s}\bm W_m\cdot\bm C_p\,\mathrm dV .
\label{eq:Kbfzi}
\end{align}
其中 \(\bm K\) 并非数学刚度矩阵；此处专用作
\textbf{海水电流模式--BFZI 端口互感耦合矩阵}。
上述基于实值模式的 \(\bm R_s,\bm L_s,\bm K\) 均为实矩阵，
其中 \(\bm R_s\) 对称正定，\(\bm L_s\) 在适用自由空间核和电流空间中对称半正定，
其物理含义分别为电阻功率与磁场能量度量。
令
\(\bm A_r(\omega)=\bm R_s+\jj\omega\bm L_s\)，
得到
\begin{equation}
\boxed{\bm A_r\bm z=-\jj\omega\bm K\bm I,\qquad
\bm z=-\jj\omega\bm A_r^{-1}\bm K\bm I.}
\label{eq:sea-current-state}
\end{equation}
因为任意非零 \(\bm v\in\mathbb C^r\) 有
\(\operatorname{Re}(\bm v^{\mathrm H}\bm A_r\bm v)
=\bm v^{\mathrm H}\bm R_s\bm v>0\)，
\(\bm A_r\) 在实频率上非奇异。实际采用适当频率与材料归一化，
避免 \(\bm R_s\) 和 \(\omega\bm L_s\) 尺度相差巨大造成数值条件恶化。

\subsection{端口散射阻抗与同源功率定理}
由自由空间 Green 核的交换对称性和相同端口/回路方向，
涡流磁矢势对第 \(p\) 个端口产生的附加链磁通为
\(\lambda_{s,p}=\sum_m[\bm K]_{mp} z_m\)。
按 \(+\jj\omega\lambda\) 的端口感应电压约定，
\(\bm v_s=\jj\omega\bm K^{\mathsf T}\bm z\)。
代入式 \eqref{eq:sea-current-state} 得到
\begin{equation}
\boxed{
\bm Z_{\rm sw}(\omega)
=\omega^2\bm K^{\mathsf T}
(\bm R_s+\jj\omega\bm L_s)^{-1}\bm K .}
\label{eq:zsw-schur}
\end{equation}
\textbf{定理（有限维互易、被动与海水功率一致性）}：
在本节所有假设成立、\(\bm R_s=\bm R_s^{\mathsf T}\succ0\)、
\(\bm L_s=\bm L_s^{\mathsf T}\succeq0\)、\(\bm K\) 实值，
且响应、端口和局部损耗均采用\textbf{同一} \(\bm W\)
和矩阵 \eqref{eq:Rs}--\eqref{eq:Kbfzi} 的条件下，
有
\begin{equation}
\bm Z_{\rm sw}=\bm Z_{\rm sw}^{\mathsf T},\quad
\operatorname{Herm}(\bm Z_{\rm sw})\succeq0,
\label{eq:sea-passivity-theorem}
\end{equation}
并对任意复端口电流 \(\bm I\) 满足
\begin{equation}
\boxed{
\frac12\bm I^{\mathrm H}\operatorname{Herm}(\bm Z_{\rm sw})\bm I
=\frac12\bm z^{\mathrm H}\bm R_s\bm z
=\int_{\Omega_s}q_s(\bm x)\,\mathrm dV,
\quad
q_s=\frac{|\bm W(\bm x)\bm z|^2}{2\sigma_s(\bm x)}.
}
\label{eq:sea-power-proof}
\end{equation}
证明：\(\bm A_r^{\mathsf T}=\bm A_r\) 导致
\((\bm A_r^{-1})^{\mathsf T}=\bm A_r^{-1}\)，故阻抗对称。
由
\[
\operatorname{Herm}(\bm A_r^{-1})
=\bm A_r^{-\mathrm H}\bm R_s\bm A_r^{-1}
\]
可得
\[
\operatorname{Herm}(\bm Z_{\rm sw})
=\omega^2\bm K^{\mathsf T}
\bm A_r^{-\mathrm H}\bm R_s\bm A_r^{-1}\bm K\succeq0.
\]
结合 \(\bm z=-\jj\omega\bm A_r^{-1}\bm K\bm I\)
得到前两个功率表达相等；将 \eqref{eq:Rs}
和 \eqref{eq:sea-reduced-current} 代入则得到体积分表达。
证毕。

该定理只证明\textbf{同一降阶模型内部}的物理一致性，
不证明 \(\bm W\) 足以逼近真实海水涡流；
积分数值误差若未一致处理也会破坏数值等式。
在 \(\omega\to0\) 且矩阵适当有界时，
\(\bm Z_{\rm sw}=\omega^2\bm K^{\mathsf T}\bm R_s^{-1}\bm K
+O(\omega^3)\)，提供低频渐近检验。

\subsection{热投影可直接取自同一电流系数}
对任意固体/海水热测试基 \(\varphi_a(\bm x)\)，
海水损耗投影为
\[
b_a^{\rm sea}
=\int_{\Omega_s}\varphi_a(\bm x)
\frac{|\bm W(\bm x)\bm z|^2}{2\sigma_s(\bm x)}\,\mathrm dV
=\frac12\bm z^{\mathrm H}\bm Q_a\bm z ,
\quad
[\bm Q_a]_{mn}=
\int_{\Omega_s}\frac{\varphi_a\bm W_m\cdot\bm W_n}{\sigma_s}\,\mathrm dV.
\]
此式建立同一涡流状态到\textbf{热空间输入}的明确接口。
若热测试基含常数分量 \(\varphi_0\equiv1\)，
其对应 \(\bm Q_0=\bm R_s\)，可直接满足总能量约束。
温度场和空间热源均可按合法场景在任意查询点继续解码，
但复杂海水外部流动应另由热/流动算子处理。
